\section*{第 \ref{ch:g12:newtons-laws} 章——牛顿运动定律}

\begin{solution}{exo:g12:newtons-laws:1}
(a) $p = 70 \times 10 = \qty{7.0e2}{kg.m/s}$。
(b) $v = \qty{13.9}{m/s}$：$p \approx \qty{1.8e4}{kg.m/s}$。
(c) $p = 0.0080 \times 800 = \qty{6.4}{kg.m/s}$。汽车 $>$ 短跑运动员 $>$
子弹：质量的五个数量级，远非速率所能弥补。
\end{solution}

\begin{solution}{exo:g12:newtons-laws:2}
$N = m(g + a_z)$：(a) $70 \times 11.01 = \qty{771}{N}$；
(b) $a_z = 0$，故 $N = mg = \qty{687}{N}$——第一定律：
速度恒定，受力相互抵消；(c) $70 \times 8.61 = \qty{603}{N}$。
\end{solution}

\begin{solution}{exo:g12:newtons-laws:3}
$a = g\sin 30^\circ = \qty{4.9}{m/s^2}$；
$v = \sqrt{2 \times 4.905 \times 3.0} = \qty{5.4}{m/s}$。质量
在 $ma$ 与 $mg\sin\alpha$ 之间约去。
\end{solution}

\begin{solution}{exo:g12:newtons-laws:4}
速度反向：$\Delta p = 0.058 \times (25 + 15) = \qty{2.3}{kg.m/s}$；
$F = 2.32/0.0050 = \qty{4.6e2}{N}$——约为 $800$ 倍的
\qty{0.57}{N} 重力，这就是接触过程中可忽略重力的原因。
\end{solution}

\begin{solution}{exo:g12:newtons-laws:5}
重力（地球作用于书，竖直向下）的反作用力为书对地球的拉力，竖直向上。
法向力（桌面作用于书，竖直向上）的反作用力为书对桌面的压力，
竖直向下。重力与法向力作用于\emph{同一}物体，且属于两种不同的相互作用——
在此二者大小相等，仅因 $\vect a = \vect 0$。
\end{solution}

\begin{solution}{exo:g12:newtons-laws:6}
$N = mg = \qty{490}{N}$。(a) $F > \mu_s N = \qty{221}{N}$。
(b) $a = (220.7 - 0.35 \times 490.5)/50 = \qty{0.98}{m/s^2}$。
(c) $F = \mu_k N = \qty{172}{N}$。
\end{solution}

\begin{solution}{exo:g12:newtons-laws:7}
(a) 当 $\tan\alpha > \mu_s$ 时开始滑动：$\alpha = \arctan 0.45 \approx
24^\circ$。(b) $a = g(\sin 30^\circ - 0.35\cos 30^\circ) =
9.81 \times 0.197 = \qty{1.9}{m/s^2}$。(c) 是：一旦滑动，
摩擦力降为 $\mu_k < \mu_s$，故
$a = g\cos\alpha\,(\tan\alpha - \mu_k) \approx \qty{0.9}{m/s^2} > 0$。
\end{solution}

\begin{solution}{exo:g12:newtons-laws:8}
$a = m_2 g/(m_1 + m_2) = 19.62/5.0 = \qty{3.9}{m/s^2}$；
$T = m_1 a = \qty{11.8}{N}$。悬挂的物块向下加速，
故其所受合力向下：$T$ 必须小于 $m_2 g$。
\end{solution}

\begin{solution}{exo:g12:newtons-laws:9}
$v_{\max} = \sqrt{\mu_s g R}$：干燥路面
$\sqrt{0.70 \times 9.81 \times 50} = \qty{18.5}{m/s} \approx
\qty{67}{km/h}$；潮湿路面 $\qty{14.0}{m/s} = \qty{50}{km/h}$。
两者 $mv^2/R$ 与 $\mu_s mg$ 均含 $m$，故约去。
\end{solution}

\begin{solution}{exo:g12:newtons-laws:10}
竖直方向：$T\cos\theta = mg$，故
$T = 2.943/\cos 25^\circ = \qty{3.2}{N}$。半径 $R = L\sin\theta =
\qty{0.507}{m}$；径向：$T\sin\theta = mv^2/R$，故
$v = \sqrt{gR\tan\theta} = \sqrt{9.81 \times 0.507 \times 0.466}
= \qty{1.5}{m/s}$.
\end{solution}

\begin{solution}{exo:g12:newtons-laws:11}
$a_z = (N - mg)/m$，其中 $mg = \qty{49.05}{N}$：$+\qty{1.0}{m/s^2}$、
$0$、$-\qty{1.0}{m/s^2}$。例如向上运行的电梯：先加速，再匀速，
最后减速。是的：若进入时已向下运动，相同读数则表示减速、匀速、
再次向下加速——体重秤读数反映的是加速度，而非
速度。
\end{solution}

\begin{solution}{exo:g12:newtons-laws:12}
(a) $0 = 60V + 4.0 \times 8.0$：$V = \qty{-0.53}{m/s}$。
(b) 球相对地面的速度为 $8.0 + V$：$0 = 60V + 4.0(8.0 + V)$，故
$V = -32/64 = \qty{-0.50}{m/s}$，球以 \qty{7.5}{m/s} 飞出。
(c) 球对她双手的推力——她投掷动作的
第三定律反作用力。
\end{solution}

\begin{solution}{exo:g12:newtons-laws:13}
无摩擦：$N\cos\theta = mg$（竖直方向）且 $N\sin\theta =
mv^2/R$（径向）；两式相除，$\tan\theta = v^2/(gR) = 400/784.8$：
$\theta \approx 27^\circ$。速度较慢时，汽车有沿弯道斜面向下滑动的趋势，
摩擦力须沿斜面向上；速度较快时则相反。
\end{solution}

\begin{solution}{exo:g12:newtons-laws:14}
(a) $250 \times 2500 = \qty{6.3e5}{kg.m/s}$，即每秒推力
为 $\qty{6.3e5}{N}$。(b) $P = \qty{4.9e5}{N}$；
$a = (6.25 - 4.905) \times 10^5 / \num{5.0e4} = \qty{2.7}{m/s^2}$.
(c) 火箭推的是自身喷出的气体（第三定律），而非空气。
\end{solution}

\begin{solution}{exo:g12:newtons-laws:15}
(a) $m_2 g > \mu_s m_1 g$：$m_2 > 0.35 \times 4.0 = \qty{1.4}{kg}$。
(b) $a = (m_2 - \mu_k m_1)\,g/(m_1 + m_2) = 1.0 \times 9.81/6.0 =
\qty{1.6}{m/s^2}$；$T = m_2(g - a) = \qty{16.4}{N}$。
(c) 令 $\mu_k = 0$，得 $a = m_2 g/(m_1 + m_2)$，
即\cref{exo:g12:newtons-laws:8}中的公式。
\end{solution}

\begin{solution}{pb:g12:newtons-laws:1}
\textbf{1.} 研究对象：车厢 $+$ 车架（$+$ 乘客）。山体
固定于地面，可视为惯性参考系，精度极高。受力：
重力 $M\vect g$（竖直向下）、轨道缆线的法向力 $\vect N$
（与轨道垂直）、牵引张力 $\vect T$（沿斜坡向上）。

\textbf{2.} 静止时由牛顿第一定律：$\sum \vect F = \vect 0$。沿
斜坡方向：$T = Mg\sin\alpha$；垂直斜坡方向：$N = Mg\cos\alpha$。

\textbf{3.} $T = 4000 \times 9.81 \times 0.500 = \qty{19.6}{kN}$；
$N = 4000 \times 9.81 \times 0.866 = \qty{34.0}{kN}$。

\textbf{4.} 空车时：$T = \qty{9.8}{kN}$。平衡中的每个力
都与 $M$ 成正比，故 $T$ 亦然。

\textbf{5.} 沿斜坡方向，$Ma = Mg\sin\alpha$：$M$ 约去，
$a = g\sin\alpha = \qty{4.9}{m/s^2}$，沿斜坡向下。

\textbf{6.} $T' = M(g\sin\alpha + a) = 4000 \times 5.505 =
\qty{22.0}{kN}$。

\textbf{7.} $a/(g\sin\alpha) = 0.60/4.905 \approx 12\%$。

\textbf{8.} $t = v/a = \qty{10}{s}$；$d = v^2/(2a) = \qty{30}{m}$。

\textbf{9.} $\vect a$ 的竖直分量：$a\sin\alpha =
\qty{0.30}{m/s^2}$，故 $N = m(g + a\sin\alpha) = 80 \times 10.11 =
\qty{809}{N}$，而 $mg = \qty{785}{N}$：重了 $3\%$。

\textbf{10.} 地面对她的静摩擦力，水平方向：
$ma\cos\alpha = 80 \times 0.52 \approx \qty{42}{N}$。身体前倾时，
地面对她的合作用力通过她的质心，
因此不会使她翻倒。

\textbf{11.} $a_T = dv/dt = 0$（匀速率）；
$a_N = v^2/R = 36/40 = \qty{0.90}{m/s^2}$，指向圆心——
在顶点处竖直向下。

\textbf{12.} 法向（取向下为正）：$Mg - N' = Mv^2/R$，于是
$N' = M(g - v^2/R) = 4000 \times 8.91 = \qty{35.6}{kN}$。

\textbf{13.} $N = m(g - v^2/R) = 80 \times 8.91 = \qty{713}{N}$：
轻了 $v^2/(gR) = 0.90/9.81 \approx 9\%$。

\textbf{14.} 当 $N' = 0$ 时，$v = \sqrt{gR} = \qty{19.8}{m/s} \approx
\qty{71}{km/h}$——约为巡航速率的三倍多，车厢
沿缆线运行时余量很大。

\textbf{15.} 由牛顿第三定律，\qty{35.6}{kN}，向下压在
缆线上（进而压在塔架上）——略小于静止时的重力，
这是由于曲线飞行而减小。

\textbf{16.} $F = \Delta p/\Delta t = 4000 \times 6.0/1.5 =
\qty{16}{kN}$，方向与运动方向相反（$a = \qty{4.0}{m/s^2}$）。

\textbf{17.} $f = ma = 80 \times 4.0 = \qty{320}{N}$；天花板可提供的最大静摩擦力为
$\mu_s mg = 0.40 \times 785 = \qty{314}{N} < \qty{320}{N}$：仅脚
会打滑——请抓住扶手。

\textbf{18.} $0 = M_0 V + m(v_{\text{相对}} + V)$：
$V = -80 \times 2.5/2080 = \qty{-0.096}{m/s}$；乘客
相对地面的速度为 $2.5 - 0.096 \approx \qty{2.4}{m/s}$。

\textbf{19.} 火箭的推力等于单位时间喷出的质量
乘以其相对火箭的喷射速率。

\textbf{20.} 遇到的最大张力出现在启动时，
$T' = M \times \qty{5.505}{N/kg}$（静止时为 $Mg\sin\alpha$，
顶点处更小）。限载：$T' \leq \qty{20}{kN}$，故
$M \leq 20000/5.505 = \qty{3.63e3}{kg}$：载荷
\qty{1.63e3}{kg}，即 $20$ 名乘客，每人 \qty{80}{kg}（而非 $25$ 名）。
标牌：\emph{20 人}；此时 $M = \qty{3.6e3}{kg}$，
$T' = \qty{19.8}{kN}$，安全系数 $100/19.8 \approx 5.1 \geq 5$。
\end{solution}
